\section{Introduction}
\label{sec:introduction}

Enterprise blockchain adoption faces a fundamental tension between transparency and
privacy. Public blockchains provide verifiable state transitions but expose sensitive
business data---pricing, supply chain relationships, production volumes---to all network
participants. Private blockchains preserve confidentiality but sacrifice the trust
guarantees that motivate blockchain adoption in the first place. Zero-knowledge validium
architectures resolve this tension by maintaining off-chain data with on-chain validity
proofs, ensuring that enterprise state transitions are cryptographically verified without
revealing the underlying data~\cite{starkware2024validium, ethereum2024rollups}.

For enterprise deployments on application-specific chains, the L1 state commitment
contract is the critical trust anchor. This contract receives ZK proofs from enterprise
nodes, verifies them, and maintains an auditable history of state root transitions. Unlike
public rollups where the state commitment contract is a shared resource, enterprise
validium requires \emph{per-enterprise state isolation}: each enterprise maintains an
independent chain of state roots, and no enterprise can interfere with another's state
history. The design of this contract directly determines the gas cost per batch submission,
the on-chain storage footprint, and the safety guarantees available to auditors and
cross-enterprise verifiers.

This paper investigates the gas-optimal storage layout for an L1 state commitment contract
on Avalanche Subnet-EVM, targeting the Basis Network enterprise ZK validium. We compare
three storage layouts---minimal (roots in mappings, metadata in events), rich (roots plus
packed metadata structs), and events-only (no per-batch persistent storage)---across gas
cost, storage footprint, and safety invariant enforcement. Critically, we decompose the gas
budget to isolate the contribution of Groth16 ZK proof verification (BN256 pairing
precompiles) from storage and logic costs.

Our key finding is that ZK pairing verification dominates the gas budget at 72\% of total
cost (205,600 of 285,756 gas), leaving only 80,156 gas for all storage operations, state
management, event emission, and access control logic. This result establishes that
\emph{integrated verification}---embedding Groth16 verification directly in the state
commitment contract rather than delegating to a separate verifier---is not merely a
design preference but an architectural requirement to remain within a 300K gas budget.
The recommended minimal layout achieves 285,756 gas per submission (first batch) with
32 bytes of persistent storage per batch, while enforcing chain continuity, gap detection,
reversal prevention, and enterprise isolation invariants with 100\% correctness across
all test scenarios.
